% ===================================================================
%  13. TAULA — Hotela. — Jatetxea. — Kafetegia.
%  Traduction 2026 d'apres texto/io, controlee sur texto/fr, texto/en
%  et texto/es. Consigne : voir texto/eu/10-taula-01.tex.
%
%  Deux notes, marquees toutes deux « (*) », et deux appels : celui du
%  bloc 01-2 (la chambre) et celui du bloc 09-1 (le menu).
%
%  DEUX CONDUCTEURS SUR LE MEME TROTTOIR, celui de l'omnibus (30) et
%  celui du fiacre (92), et deux commissionnaires, le chasseur (19) et
%  l'homme de peine (39) : quatre metiers voisins qu'objekti.py releve
%  par leur nom et qui doivent donc porter quatre mots distincts.
% ===================================================================

%%K t13-noto-1 noto
\VUnotes{143.2mm}{%
\textit{\VUgras{(*) Logelaren xehetasunak ikus 6. taulan.}}%
}

%%K t13-apar-1 apar
\VUsaut{3.0mm}
\VUtitre{15.0pt}{60}{13. TAULA}
\VUsaut{1.6mm}
\VUpk{10.2pt}{40}{Hotela. --- Jatetxea.}
\VUsaut{2.0mm}
\VUpk{10.2pt}{40}{Kafetegia.}
\VUsaut{3.4mm}

%%K t13-01-1 p
1. --- Atzo arratsean Royanera iritsita, \textit{Ozeanoaren Hotelean} hartu nuen
ostatu, adiskide batek gomendatuta.

%%K t13-01-2 p
\VUgras{Gela} eder bat \textsuperscript{(2)} eman zidaten bigarren
\VUgras{solairuan} \textsuperscript{(1)}, non ederki lo egin bainuen (*).

%%K t13-02-1 p
2. --- Gaur goizean, nire gelatik atera nintzenean --- non \VUgras{neskameak}
\textsuperscript{(3)} gosaria ekarri baitzidan ---, \VUgras{erretzeko gelara}
\textsuperscript{(5)} sartu nintzen, zeina \VUgras{irakurgelatzat}
\textsuperscript{(4)} ere erabiltzen baita, \VUgras{egunkariak}
\textsuperscript{(6)} irakurtzeko, \VUgras{zigarreta} bat
\textsuperscript{(8)} erretzen nuen bitartean. Irakurri ondoren,
\VUgras{igogailuan} \textsuperscript{(9)} jaitsi nintzen eta hirian barrena
paseatzera joan.

%%K t13-03-1 p
3. --- Hoteleko \VUgras{bulego-langileari} \textsuperscript{(12)} galdetzea ahaztu
zitzaidanez zein ordutan zerbitzatzen duten \VUgras{mahai komunean}
\textsuperscript{(11)}, goiz itzuli nintzen eta hura baliatu nuen hoteleko
\VUgras{bainugelan} \textsuperscript{(10)} bainatzera joateko. Gero
\VUgras{kutxara} \textsuperscript{(15)} joan nintzen adiskide bati deitzeko. Baina
\VUgras{telefonoa} \textsuperscript{(13)} okupatuta zegoen; nire estutasuna
ikusirik, \VUgras{kutxazainak} \textsuperscript{(14)}, zeinak \VUgras{faktura}
bat \textsuperscript{(17)} sartzen baitzuen \VUgras{bidaiarien liburuan}
\textsuperscript{(16)}, \VUgras{atezainari} \textsuperscript{(18)} eskatu zion
\VUgras{mandatari} bat \textsuperscript{(19)} bidaltzeko nire mandatua egitera bere
\VUgras{bizikletan} \textsuperscript{(20)}.

%%K t13-04-1 p
4. --- Eskerrak eman nizkion eta \VUgras{zamaketariari} \textsuperscript{(24)}
eskupekoa ematera atera nintzen (gauza guztietarako gizona), zeinak geltokitik
bere \VUgras{esku-organ} \textsuperscript{(25)} ekarri baitzituen nire
\VUgras{kapela-kartoia} \textsuperscript{(26)}, nire \VUgras{maleta}
\textsuperscript{(27)} eta nire \VUgras{bidaia-zakua} \textsuperscript{(28)}.
\VUgras{Postariak} \textsuperscript{(21)}, une hartan bertan iritsi zenak,
\VUgras{gutun} bat \textsuperscript{(22)} eman zidan.

%%K t13-05-1 p
5. --- Pixka bat gehiago geratu nintzen atalasean, \VUgras{postontziaren}
\textsuperscript{(23)} ondoan, kaleko mugimenduari begira. \VUgras{Gidariak}
\textsuperscript{(30)}, \VUgras{omnibusarenak} \textsuperscript{(29)},
\VUgras{kutxa} bat \textsuperscript{(31)} kargatzen zuen bere gurdian,
\VUgras{bidaiari} batena \textsuperscript{(32)}, zeinari \VUgras{hoteleko
jabeak} \textsuperscript{(33)} agur egiten baitzion. \VUgras{Telegrafista}
gazte batek \textsuperscript{(35)} \VUgras{despatxu} bat \textsuperscript{(34)}
zeraman, batere presarik gabe, hoteleko bezero bati; \VUgras{turista} batek
\textsuperscript{(36)} bere \VUgras{argazki-makina} \textsuperscript{(38)}
bizkarrean, bere \VUgras{gidaliburua} \textsuperscript{(37)} kontsultatzen zuen.

%%K t13-tit-1 sub
\VUsaut{4.0mm}
\VUpk{10.2pt}{40}{Gosariaren ordua.}
\VUsaut{3.0mm}

%%K t13-06-1 p
6. --- Burdinsarearen aurrean \VUgras{mandazain} axolagabe bat
\textsuperscript{(39)} lo-kuluxka egiten ari zen bere \VUgras{kakoen}
\textsuperscript{(40)} eta bere \VUgras{zapata-garbitzailearen kutxaren}
\textsuperscript{(41)} artean, eta \VUgras{ikuzle} gazte batek
\textsuperscript{(42)}, arropa-\VUgras{saski} bat \textsuperscript{(43)}
zeramanak, barre egiten zuen atean zegoen \VUgras{mahai-nagusiaren}
\textsuperscript{(44)} itxura serioaz.

%%K t13-07-1 p
7. --- \VUgras{Sukaldariaren} \textsuperscript{(45)} itxurak --- zeina bere \VUgras{sukaldetik}
\textsuperscript{(47)} atera baitzen, \VUgras{sotokotik}
\textsuperscript{(48)}, \VUgras{sukalde-morroi} bat
\textsuperscript{(46)} mandatu batera bidaltzeko --- gogorarazi zidan gosariaren
ordua hurbil zela, eta jatetxera joan nintzen, patioaren beste aldera, non
\VUgras{automobilista} batek \textsuperscript{(50)} bere \VUgras{automobila}
\textsuperscript{(49)} sartzen baitzuen, eta \VUgras{mekanikari} batek
\textsuperscript{(52)} \VUgras{txirrindulariaren} \textsuperscript{(53)}
argibideen arabera \VUgras{gasolinazko triziklo} bat \textsuperscript{(51)}
konpontzen baitzuen, istripua izan berri zuena.

%%K t13-08-1 p
8. --- \VUgras{Mutil} gazte bati \textsuperscript{(54)}, zeinak \VUgras{garagardo
pitxerrak} \textsuperscript{(55)} baitzeramatzan, galdetu nion ea mahai komuna
berandaren azpian zegoen, eta baietz esan zidanean, mahaietako batean eseri
nintzen eta gosari bikaina eman zidaten.

%%K t13-noto-2 noto
\VUnotes{143.2mm}{%
\textit{\VUgras{(*) Hiztegian emandako jaki-zerrendaren laguntzaz, irakasleak
erraz alda dezake behean emandako menua.}}%
}

%%K t13-tit-2 sub
\VUsaut{4.0mm}
\VUpk{10.2pt}{40}{Gosari bikaina.}
\VUsaut{3.0mm}

%%K t13-09-1 p
9. --- Zerbitzariak gosariaren menua (*) eta ardo-karta ekarri zizkidan. Aukeratu
eta eskatu nituen \VUgras{aurreko platertzat} \VUgras{izkirak}
\VUgras{gurinarekin}, eta, lehen platertzat, \VUgras{Caen erako tripakiak}.
\VUgras{Arrainik} ez nuen hartu, baina arkume-\VUgras{hanka} zati bat eskatu nuen,
eta, \VUgras{barazkitzat, ziazerbak} esnegainarekin. Gero, \VUgras{plater
gozoetako} bat ere ez zitzaidanez gustatu, postre desberdinak ekartzeko agindu
nuen: \VUgras{gazta}, \VUgras{fruta} eta \VUgras{gaileta}. Saint-Émilion
\VUgras{ardo} bikain bat gehitu nuen, eta nire gosariarekin oso pozik, mahaitik
altxatu nintzen \VUgras{zerbitzariari} \textsuperscript{(57)} \VUgras{eskupeko}
eder bat \textsuperscript{(59)} eman ondoren.

%%K t13-10-1 p
10. --- Alde egitera nindoala ikusirik, \VUgras{zuzendariak}
\textsuperscript{(60)} balkoira ateratzeko proposatu zidan \VUgras{ozeanoaren}
\textsuperscript{(62)} panorama bikaina miresteko. Urrunean arrantzako
\VUgras{txalupa} txikiak \textsuperscript{(63)}, \VUgras{belaontziak}
\textsuperscript{(64)} eta \VUgras{paketebote} erraldoi bat
\textsuperscript{(65)} bereizten ziren, zeinaren silueta beltza \VUgras{hodei}
zurien gainean \textsuperscript{(67)} ebakitzen baitzen \VUgras{zeruertzean}
\textsuperscript{(66)}. Haizetxo arin batek \VUgras{bandera}
\textsuperscript{(70)} astintzen zuen, \VUgras{iragarkiaren}
\textsuperscript{(69)} gainean jarria, \VUgras{atondokoaren}
\textsuperscript{(68)} gainean.

%%K t13-tit-3 sub
\VUsaut{3.0mm}
\VUpk{10.2pt}{40}{Denbora-pasak.}
\VUsaut{3.0mm}

%%K t13-11-1 p
11. --- Ikuskizuna hain zen interesgarria, non erabaki bainuen hurrengo egunean
kafetegiaren terrazara igotzea \VUgras{antzoki-betaurrekoekin}
\textsuperscript{(71)} edo \VUgras{teleskopioarekin} \textsuperscript{(72)}.

%%K t13-12-1 p
12. --- Balkoitik alde eginda, hoteleko kafetegira joan nintzen \VUgras{zerbait}
\textsuperscript{(73)} edatera eta \VUgras{billar} \textsuperscript{(75)} partida
bat jokatzera. Zuzenean zeharkatu nuen kafetegiko aretoa, non bezero askok
\VUgras{xakera} \textsuperscript{(78)}, \VUgras{dametara}
\textsuperscript{(79)}, \VUgras{dominora} \textsuperscript{(80)},
\VUgras{triktrakera} \textsuperscript{(81)} edo \VUgras{kartetara}
\textsuperscript{(82)} jokatzen baitzuten, eta zuzenean igo nintzen
\VUgras{billar-gelara} \textsuperscript{(74)}.

%%K t13-13-1 p
13. --- Partida amaituta, beheko solairura jaitsi nintzen eta zerbitzariari
\VUgras{gutun-azal} bat \textsuperscript{(84)}, \VUgras{papera}
\textsuperscript{(83)}, \VUgras{seilu} bat \textsuperscript{(85)},
\VUgras{luma} bat \textsuperscript{(87)} eta \VUgras{tinta}
\textsuperscript{(86)} eskatu nizkion gutun bat idazteko.

%%K t13-13-2 p
Gero \VUgras{kotxezain} bat \textsuperscript{(92)} deitu nuen, zeinaren
\VUgras{kotxea} \textsuperscript{(88)} kafetegiaren aurrean gelditu baitzen.
Orduko edo bidaiako prezioa galdetu nion, eta, prezioan ados jarri ginenean,
\VUgras{atetxoa} \textsuperscript{(89)} ireki zidan eta barruan jarri ninduen.
Bere \VUgras{pesketara} \textsuperscript{(91)} berriro igo zen, zaldiak bizkor
abiarazi zituen bere \VUgras{zigor} txikiaz \textsuperscript{(93)}, eta ibilaldi
zoragarri batera abiatu ginen.
\VUsaut{6.0mm}
\VUfilet{22mm}
